\section{材料}

実験で使用した材料を以下に示す。 
  \subsection{ITO付きガラス基板（ITO）}
    \begin{tabular}{ll}
        製造&ジオマテック株式会社\\
        品番&0031\\
        サイズ&25\,$\times$\,25\,$\times$\,0.7 \si{mm}\\
    \end{tabular}
    \\
    本研究では、spiro-OMeTADまたは中間層を堆積させる透明電極基板として用いる。ITOはIndium Tin Oxideの略称で、日本語では酸化インジウムスズと呼ばれる。ITO付きガラス基板は、洗浄後、ソックスレー抽出機を用いて2-プロパノールに浸漬し、保管する。\\
    

   \subsection{spiro-OMeTAD}
    \begin{tabular}{ll}
    化学名&2,2',7,7'-Tetrakis(N,N-di-p-methoxyphenylamino)-9,9'-spirobifluorene\\
    製造&東京化成株式会社\\
    品番&T3672\\
    化学式&\ce{C81H68N4O8} \\
    分子量&1225.43\\
    純度&\SI{99}{\%}\\
    %CAS    &  207739-72-8\\
    \end{tabular}
    \\
 spiro-OMeTADは、ペロブスカイト太陽電池のホール輸送材料としてよく使用されているが、本研究ではEC材料として使用する。spiro-OMeTADは、1価、2価、4価時にそれぞれ赤色、青色、紫色になることが分かっている。\cite{山本雄樹}$\,$デシケーター内で保管する。\\
 spiro-OMeTADの構造式を図\ref{spiro-OMeTAD}に示す。
 \begin{figure}[H]
	\centering
	\includegraphics[width=8cm]{Spiro-OMeTAD.png}
	\caption{spiro-OMeTADの構造式}
	\label{spiro-OMeTAD}
\end{figure}


   \subsection{カフェ酸}
    \begin{tabular}{ll}
    製造&東京化成株式会社\\
    品番&C0002\\
    化学式&\ce{C9H8O4}\\
    分子量&180.16\\
    純度&\SI{98.0}{\%}以上\\
    \end{tabular}
    \\
    本研究で中間層に用いる材料である。デシケーター内で保管する。\\
     カフェ酸の構造式を図\ref{Cafe}に示す。
 \begin{figure}[H]
	\centering
	\includegraphics[width=7cm]{Cafe acid.png}
	\caption{caffeic acidの構造式}
	\label{Cafe}
\end{figure}

\subsection{2PACz}
    \begin{tabular}{ll}
    製造&東京化成株式会社\\
    品番&C3663\\
    化学式&\ce{C14H14NO3P}\\
    分子量&275.24\\
    純度&\SI{98.0}{\%}以上\\
    \end{tabular}
    \\
    本研究で中間層に用いる材料である。デシケーター内で保管する。\\
     2PACzの構造式を図\ref{Cz}に示す。
 \begin{figure}[H]
	\centering
	\includegraphics[width=7cm]{2PACz.png}
	\caption{2PACzの構造式}
	\label{Cz}
\end{figure}
    
    \subsection{Br-2PACz}
    \begin{tabular}{ll}
    製造&東京化成株式会社\\
    品番&C3663\\
    化学式&\ce{C14H12Br2NO3P}\\
    分子量&433.04\\
    純度&\SI{99.0}{\%}以上\\
    \end{tabular}
    \\
    本研究で中間層に用いる材料である。デシケーター内で保管する。\\
     Br-2PACzの構造式を図\ref{Br}に示す。
 \begin{figure}[H]
	\centering
	\includegraphics[width=7cm]{Br-2PACz.png}
	\caption{Br-2PACzの構造式}
	\label{Br}
\end{figure}

    \subsection{Me-2PACz}
    \begin{tabular}{ll}
    製造&東京化成株式会社\\
    品番&M3477\\
    化学式&\ce{C16H18NO3P}\\
    分子量&303.30\\
    純度&\SI{99.0}{\%}以上\\
    \end{tabular}
    \\
    本研究で中間層に用いる材料である。デシケーター内で保管する。\\
     Me-2PACzの構造式を図\ref{Me}に示す。
 \begin{figure}[H]
	\centering
	\includegraphics[width=7cm]{Me-2PACz.png}
	\caption{Me-2PACzの構造式}
	\label{Me}
\end{figure}

\subsection{過塩素酸ナトリウム}
    \begin{tabular}{ll}
    製造&Sigma-Aldrich\\
    品番&410241\\
    化学式&\ce{NaClO4}\\
    分子量&122.44\\
    純度&\SI{98.0}{\%}以上
    \end{tabular}
    \\
    本研究では、電気化学測定時の電解質として使用した。デシケーター内で保管した。\\
    過塩素酸ナトリウムの構造式を図\ref{Na}に示す。
 \begin{figure}[H]
	\centering
	\includegraphics[width=7cm]{NaClO4.png}
	\caption{過塩素酸ナトリウムの構造式}
	\label{Na}
\end{figure}


    \subsection{クロロベンゼン}
    \begin{tabular}{ll}
    製造&富士フイルム和光純薬株式会社\\
    品番&032-07986\\
    化学式&\ce{C6H5Cl}\\
    分子量&112.56\\
    純度&\SI{99.0}{\%}以上\\
    %CAS&443922-06-3\\
    \end{tabular}
    \\
    EC薄膜を成膜する際に用いるEC材料を溶かす溶媒として使用する。アルミパックの中に入れ暗所で保管する。\\

　　 \subsection{エタノール}
	 \begin{tabular}{ll}
     製造&富士フイルム和光純薬株式会社\\
     品番&057-00456\\
     化学式&\ce{C2H5OH}\\
     分子量&46.07\\
     純度&\SI{99.5}{\%}以上\\
     \end{tabular}
     \\
     中間層を成膜する際に用いる材料を溶かす溶媒として使用する。暗所で保管する。\\

    \subsection{脱イオン水}
    \noindent 電解液の溶媒と、基板の洗浄に脱イオン水を用いる。脱イオン水は、純水装置（オルガノ株式会社、ピュアライト RPA-0015-001）を用いて、水道水から作製した。\\

    \subsection{2-プロパノール}
    \begin{tabular}{ll}
    製造&富士フイルム和光純薬株式会社\\
    化学式&\ce{CH3CH(OH)CH3}\\
    分子量&60.10\\
    \end{tabular}
    \\
    \noindent 本研究ではITO付きガラス基板を洗浄する際の溶媒として利用した．容器は遮光し，なるべく涼しい場所で保管した．\\